\subsection{Area utente}

Nella parte in alto a destra della schermata trovi il menu utente, da cui puoi accedere rapidamente alle principali funzionalit\`a del tuo account.

% TODO: inserire figura posizione menu utente

Cliccando sul tuo nome, si apre un menu a tendina con diverse opzioni.

% TODO: inserire figura menu utente aperto

Da qui puoi:
\begin{itemize}[itemsep=5pt, parsep=5pt, label=$\scriptstyle\bullet$]
	\item Profilo, per visualizzare e gestire i tuoi dati
	\item Chat, per tornare alla schermata principale
	\item Storico, per consultare i tuoi ordini
	\item Logout, per uscire dall'applicazione
\end{itemize}

% TODO: inserire figura opzioni del menu utente

Questo menu ti permette di muoverti facilmente tra le diverse sezioni senza interrompere il tuo lavoro.

\subsubsection{Profilo utente}

Nella sezione Profilo puoi visualizzare le informazioni principali del tuo account, come:
\begin{itemize}[itemsep=5pt, parsep=5pt, label=$\scriptstyle\bullet$]
	\item Username
	\item Email
\end{itemize}

% TODO: inserire figura sezione profilo utente

Da qui puoi anche gestire alcune impostazioni importanti.

\subsubsubsection{Modifica password}

Se vuoi cambiare la password, clicca su \vr{Reimposta password}.

% TODO: inserire figura pulsante reimposta password

Ti verr\`a chiesto di inserire:
\begin{itemize}[itemsep=5pt, parsep=5pt, label=$\scriptstyle\bullet$]
	\item La password attuale
	\item La nuova password
	\item La conferma della nuova password
\end{itemize}

% TODO: inserire figura form modifica password

Dopo aver compilato i campi, puoi confermare per aggiornare la password.

\subsubsubsection{Eliminazione account}

Se desideri eliminare il tuo account, puoi utilizzare il pulsante \vr{Cancella account}.

% TODO: inserire figura pulsante cancella account

Prima di procedere, il sistema ti mostrer\`a una finestra di conferma per avvisarti che l'operazione \`e definitiva.

% TODO: inserire figura conferma eliminazione account

Puoi quindi scegliere se:
\begin{itemize}[itemsep=5pt, parsep=5pt, label=$\scriptstyle\bullet$]
	\item Confermare, eliminando definitivamente l'account
	\item Annullare, per tornare indietro senza modifiche
\end{itemize}

\subsubsection{Storico ordini}

Nella sezione Storico puoi visualizzare tutti gli ordini che hai effettuato.

% TODO: inserire figura sezione storico ordini

Ogni ordine mostra alcune informazioni utili, come:
\begin{itemize}[itemsep=5pt, parsep=5pt, label=$\scriptstyle\bullet$]
	\item Codice ordine
	\item Data
	\item Numero di prodotti
	\item Accesso al dettaglio dell'ordine
\end{itemize}

% TODO: inserire figura elenco ordini nello storico

Selezionando un ordine, puoi vedere nel dettaglio tutti i prodotti inclusi.\\
Inoltre, puoi duplicare un ordine, cos\`i da riutilizzarlo rapidamente senza dover reinserire tutti i prodotti.

\subsubsubsection{Filtri e navigazione storico}

Per aiutarti a trovare pi\`u facilmente gli ordini, sono disponibili alcune funzioni:

\begin{itemize}[itemsep=5pt, parsep=5pt, label=$\scriptstyle\bullet$]
	\item \textbf{Filtro per data}: puoi selezionare un intervallo di date per visualizzare solo gli ordini effettuati in quel periodo.
\end{itemize}

% TODO: inserire figura filtro per data nello storico

\begin{itemize}[itemsep=5pt, parsep=5pt, label=$\scriptstyle\bullet$]
	\item \textbf{Paginazione}: se hai molti ordini, puoi scorrere le pagine utilizzando i pulsanti:
\end{itemize}
\begin{itemize}[itemsep=5pt, parsep=5pt, label=$\scriptstyle\bullet$]
	\item \vr{Precedente}
	\item \vr{Successivo}
\end{itemize}

% TODO: inserire figura paginazione storico

\subsubsection{Logout}

Per uscire dall'applicazione, seleziona \vr{Logout} dal menu utente.

% TODO: inserire figura opzione logout

Verrai disconnesso e riportato alla schermata di login.
